% GENERATED FILE. Edit canonical lessons, then run npm run generate:books.
% canonical-source-hash: fnv1a64:40033eb742ac09f8
% canonical-lessons: TE-C34-tiisuko, TE-C34-adugu, TE-C34-sahayam-ceyu, TE-C34-ishtam

\chapter{Four Verbs Between People}
\label{ch:te-verbs-between-people}
\hlchaptermodality{34}

\begin{chapteropening}
\textbf{By the end of this chapter:} \emph{I can say I take, I ask and I help in Telugu, say \te{నాకు} \te{తెలుగు} \te{ఇష్టం} knowing there is no verb in it, place it beside the three other dative sentences I have built, and name the branch of Dravidian Telugu actually belongs to.}

Say \te{నాకు} \te{తెలుగు} \te{ఇష్టం} and unpack it as \textquotedblleft{}to me, Telugu, a liking\textquotedblright{} with no verb at all; line it up against \te{నాకు} \te{తెలుగు} \te{వచ్చు}, \te{నాకు} \te{తెలుసు} and \te{నాకు} \te{అర్థమైంది}; choose between the native \te{నచ్చు} and \te{ఇష్టపడు}, whose \te{పడు} is the falling of \te{వర్షం} \te{పడుతోంది}; and trace \te{ఇష్టం} to the Indo-European root behind English ask.
\end{chapteropening}

\section[tīsukō]{\te{తీసుకో} (tīsukō) — \textquotedblleft{}to take,\textquotedblright{} the third verb you do for yourself}

\label{lesson:TE-C34-tiisuko}



\begin{quote}
The last four verbs happened inside your head; these four happen between you and somebody else. Say \emph{caduvutānu} and \emph{rāstānu}, and name what \emph{vrāyu}'s root first meant.
\end{quote}

\subsection*{You'll want to know: \te{తీసుకో}}
\begin{quote}
\textbf{\te{తీసుకో}} — \emph{tīsukō} — \textbf{to take}
\end{quote}

Underneath it is \textbf{\te{తీయు}} (\emph{tīyu}), \textquotedblleft{}to pull out,\textquotedblright{} and then the \textbf{-\te{కో}} you have met twice: pull it out, \textbf{and keep it}.

\textbf{\te{తీసుకుంటాను}} (\emph{tīsukuṇṭānu}), \textquotedblleft{}I take\textquotedblright{}; \textbf{\te{తీసుకున్నాను}} (\emph{tīsukunnānu}), \textquotedblleft{}I took\textquotedblright{} — \emph{tiṇṭānu}'s ending on a different stem. Two whole sentences, from Chapter 15:

\begin{quote}
\textbf{\te{నేను} \te{నీళ్ళు} \te{తీసుకుంటాను}.} — \emph{nēnu nīḷḷu tīsukuṇṭānu} — \textquotedblleft{}I take water.\textquotedblright{}
\end{quote}

>

\begin{quote}
\textbf{\te{నేను} \te{బియ్యం} \te{తీసుకుంటాను}.} — \emph{nēnu biyyaṁ tīsukuṇṭānu} — \textquotedblleft{}I take rice.\textquotedblright{}
\end{quote}

\subsection*{The letters in this word}
\textbf{\te{తీ}} is the dental \textbf{\te{త}} (\emph{ta}) with the long-\emph{ī} sign — \emph{tinu}'s letter, held longer. \textbf{\te{సు}} is \textbf{\te{స}} (\emph{sa}) plus \emph{u}; \textbf{\te{కో}} is \emph{anukō}'s \emph{kō}.

\begin{grammarlens}[title={Grammar Lens: three verbs, one ending}]
\noindent
\begin{tabularx}{\linewidth}{@{}>{\raggedright\arraybackslash}X>{\raggedright\arraybackslash}X@{}}
\toprule
\textbf{the verb} & \textbf{what the -\te{కో} adds} \\
\midrule
\emph{anukō} & say it — \textbf{to yourself} \\
\emph{arthaṁ cēsukō} & make the meaning — \textbf{yours} \\
\emph{tīsukō} & pull it out — \textbf{for yourself} \\
\bottomrule
\end{tabularx}

The respectful \textbf{-\te{అండి}} goes on this one too: \textbf{\te{తీసుకోండి}} (\emph{tīsukōṇḍi}), \textquotedblleft{}please take\textquotedblright{} — what a Telugu host says while pushing a dish towards you.
\end{grammarlens}

\begin{cousinweb}
\textbf{\te{కొను}} (\emph{konu}) is still a verb on its own: \textquotedblleft{}to take, to get, to buy.\textquotedblright{} Its cousins are exact — Tamil \textbf{\ta{கொள்}} (\emph{koḷ}) and Malayalam \textbf{\ml{കൊള്ളുക}} (\emph{koḷḷuka}), both \textquotedblleft{}to hold, to take.\textquotedblright{}

What those two did with it is the good part: they welded it onto the backs of other verbs as an ending, with the very meaning Telugu gives it. So the shape you learned two lessons ago is not a Telugu quirk. It is Dravidian machinery, and \emph{tīsukō} is the verb it started from.
\end{cousinweb}

\subsection*{Your turn}
Say these aloud:

\begin{itemize}
\raggedright
  \item \textquotedblleft{}tīyu … kō\textquotedblright{} $\to$ \textquotedblleft{}tīsukō\textquotedblright{}
  \item \textquotedblleft{}nēnu nīḷḷu tīsukuṇṭānu\textquotedblright{} — then with \textquotedblleft{}biyyaṁ\textquotedblright{}, rice
  \item the three with the ending — \textquotedblleft{}anukō … arthaṁ cēsukō … tīsukō\textquotedblright{}
  \item what a host says — \textquotedblleft{}tīsukōṇḍi\textquotedblright{}
  \item the cousins — \textquotedblleft{}konu … koḷ … koḷḷuka\textquotedblright{}
\end{itemize}

\begin{tcolorbox}[breakable,colback=teal!4,colframe=teal!35!black,title={Before you move on}]
Take \emph{tīsukō} apart, say \textquotedblleft{}I take water,\textquotedblright{} and name the other two verbs with that ending. (\emph{\textbf{\te{తీయు}}} plus \textbf{-\te{కో}}; \emph{nēnu nīḷḷu tīsukuṇṭānu}; \emph{anukō} and \emph{arthaṁ cēsukō}.) Which verb did the ending come from, and which sisters have it? (\emph{\textbf{\te{కొను}}}; \textbf{Tamil} \emph{koḷ}, \textbf{Malayalam} \emph{koḷḷuka} — both use it as an ending too.) How do you offer someone food? (\emph{\textbf{\te{తీసుకోండి}}}.)
\end{tcolorbox}
\section[aḍugu]{\te{అడుగు} (aḍugu) — \textquotedblleft{}to ask,\textquotedblright{} and the branch Telugu really sits on}

\label{lesson:TE-C34-adugu}



\begin{quote}
\emph{Tīsukuṇṭānu} — and its ending belonged to the whole family. This next verb belongs to almost nobody.
\end{quote}

\subsection*{You'll want to know: \te{అడుగు}}
\begin{quote}
\textbf{\te{అడుగు}} — \emph{aḍugu} — \textbf{to ask; to request}
\end{quote}

\textbf{\te{అడుగుతాను}} (\emph{aḍugutānu}), \textquotedblleft{}I ask\textquotedblright{}; \textbf{\te{అడిగాను}} (\emph{aḍigānu}), \textquotedblleft{}I asked\textquotedblright{}; \textbf{\te{అడగండి}} (\emph{aḍagaṇḍi}), \textquotedblleft{}please ask.\textquotedblright{} \emph{aḍu-}, \emph{aḍi-}, \emph{aḍa-}: the middle vowel wanders as \emph{caduvu}'s did, and the polite \textbf{-\te{అండి}} rides on unbothered.

One warning: Telugu has a \textbf{second} \te{అడుగు}, \textquotedblleft{}a foot, a step\textquotedblright{} — same spelling, different root, nothing to do with asking.

\subsection*{The letters in this word}
\textbf{\te{అ}} is the standing vowel \emph{a}, as in \emph{anukō}; \textbf{\te{డు}} is the retroflex \textbf{\te{డ}} (\emph{ḍa}) plus \emph{u}, the \emph{ḍu} ending \emph{cūḍu}; \textbf{\te{గు}} is \textbf{\te{గ}} (\emph{ga}) plus \emph{u}.

\begin{cousinweb}
The four literary sisters share no ask-verb: Tamil \textbf{\ta{கேள்}} (\emph{kēḷ}) and Kannada \textbf{\ka{ಕೇಳು}} (\emph{kēḷu}), which also mean \textquotedblleft{}hear\textquotedblright{}; Malayalam \textbf{\ml{ചോദിക്കുക}} (\emph{cōdikkuka}); and Telugu \textbf{\te{అడుగు}}, which keeps hearing separate in \textbf{\te{విను}} (\emph{vinu}). The Dravidian dictionary's firmest match for \emph{aḍugu} is in none of them: it is \textbf{Parji} \emph{aḍ-}, a small language of central India.

That is the clue to something worth naming outright. \textbf{Telugu is not in the same branch as Tamil, Kannada and Malayalam.} Those three are \textbf{South Dravidian}; Telugu heads \textbf{South-Central Dravidian}, with Gondi, Konda, Kui, Kuvi, Pengo and Manda. So the pattern was no run of accidents: \emph{lēdu} against the sisters' \emph{il-}, \emph{uṇḍu} against their \emph{iru}, \emph{enimidi} and \emph{tommidi} where \emph{eṭṭu} and \emph{oṉpatu} will not reach.
\end{cousinweb}

\begin{cousinweb}
For the noun Telugu reaches for Sanskrit: \textbf{\te{ప్రశ్న}} (\emph{praśna}), \textquotedblleft{}a question,\textquotedblright{} from \textbf{\dv{प्रच्छ्}} (\emph{prach}) — an Indo-European root, and the root of Latin \emph{precārī}, which English took as \textbf{pray}. Native verb, borrowed noun: \emph{tiṇḍi} and \emph{bhōjanaṁ}, a third time.
\end{cousinweb}

\subsection*{Your turn}
Say these aloud:

\begin{itemize}
\raggedright
  \item the three vowels — \textquotedblleft{}aḍugutānu … aḍigānu … aḍagaṇḍi\textquotedblright{}
  \item the same wobble — \textquotedblleft{}caduvutānu … cadivānu … cadavaṇḍi\textquotedblright{}
  \item four sisters, four verbs — \textquotedblleft{}aḍugu … kēḷ … kēḷu … cōdikkuka\textquotedblright{}
  \item Telugu's own branch — \textquotedblleft{}Gondi, Konda, Kui\textquotedblright{}
  \item what it explains — \textquotedblleft{}lēdu … uṇḍu … enimidi, tommidi\textquotedblright{}
  \item the borrowed noun and its cousin — \textquotedblleft{}praśna … pray\textquotedblright{}
\end{itemize}

\begin{tcolorbox}[breakable,colback=teal!4,colframe=teal!35!black,title={Before you move on}]
Say \textquotedblleft{}I ask,\textquotedblright{} \textquotedblleft{}I asked,\textquotedblright{} \textquotedblleft{}please ask,\textquotedblright{} and name the other verb that wanders its middle vowel. (\emph{Aḍugutānu, aḍigānu, aḍagaṇḍi}; \textbf{\te{చదువు}}.) Do the four literary sisters share an ask-verb? (\textbf{No} — \emph{aḍugu}, \emph{kēḷ}, \emph{kēḷu}, \emph{cōdikkuka}.) Which branch is Telugu on, who is on it, and name two things that explains. (\textbf{South-Central Dravidian}, with Gondi, Konda and Kui; \emph{\textbf{lēdu}} against \emph{il-}, \emph{\textbf{uṇḍu}} against \emph{iru}, \emph{\textbf{enimidi}} and \emph{\textbf{tommidi}}.) Where is \emph{praśna} from? (\textbf{Sanskrit}, from the root behind English \textbf{pray}.)
\end{tcolorbox}
\section[sahāyaṁ cēyu]{\te{సహాయం} \te{చేయు} (sahāyaṁ cēyu) — \textquotedblleft{}to help,\textquotedblright{} which is a going-together}

\label{lesson:TE-C34-sahayam-ceyu}



\begin{quote}
\emph{Aḍagaṇḍi} — please ask. Say it, and name where \emph{praśna} came from. The verb that answers is one Chapter 5 walked past.
\end{quote}

\subsection*{You'll want to know: \te{సహాయం} \te{చేయు}}
\begin{quote}
\textbf{\te{సహాయం} \te{చేయు}} — \emph{sahāyaṁ cēyu} — \textbf{to help}
\end{quote}

Chapter 5 listed it once, as an example of what \emph{cēyu} does with a noun in front. Take it properly now.

\begin{quote}
\textbf{\te{నేను} \te{మీకు} \te{సహాయం} \te{చేస్తాను}.} — \emph{nēnu mīku sahāyaṁ cēstānu} — \textquotedblleft{}I help you.\textquotedblright{}
\end{quote}

The person helped goes in the \textbf{dative}: \emph{mīku}, \textquotedblleft{}to you,\textquotedblright{} Chapter 6's \textbf{-\te{కు}} at ordinary work. Politely: \textbf{\te{సహాయం} \te{చేయండి}} (\emph{sahāyaṁ cēyaṇḍi}). Speech shortens the noun to \textbf{\te{సాయం}} (\emph{sāyaṁ}).

\subsection*{The letters in this word}
\textbf{\te{స}} is \emph{sa}; \textbf{\te{హా}} is \textbf{\te{హ}} (\emph{ha}) with the long-\emph{ā} sign; \textbf{\te{యం}} is \textbf{\te{య}} (\emph{ya}) closed by the \textbf{anusvāra}.

\begin{cousinweb}
\textbf{\te{సహాయం}} is Sanskrit \textbf{\dv{सहाय}} (\emph{sahāya}), and it comes apart cleanly: \textbf{\dv{सह}} (\emph{saha}), \textquotedblleft{}with,\textquotedblright{} plus \textbf{\dv{आय}} (\emph{āya}), \textquotedblleft{}a going.\textquotedblright{} A helper is \textbf{one who goes with you}. Both halves are Indo-European — \emph{saha} is the \textquotedblleft{}one, together\textquotedblright{} root behind \textbf{same} and \textbf{similar}, and \emph{āya} is the go-root Latin carries as \emph{īre}, which English took whole in \textbf{exit} and \textbf{transit}.

The build is Chapter 8's: a borrowed noun in front of the native \textbf{\te{చేయు}}, as \textbf{\te{దయచేసి}} was made from \emph{daya}. Kannada and Malayalam borrowed \emph{sahāya} too; \textbf{Tamil did not}, keeping its own \textbf{\ta{உதவி}} (\emph{udavi}) — yet for \emph{please} all four reached for \emph{daya}. Borrowing happens \textbf{word by word}.
\end{cousinweb}

\begin{culture}
Dictionaries record \textbf{\te{సహాయం}} and \textbf{\te{సాయం}} side by side, \emph{sāyaṁ} the shorter and more casual. What no source consulted here states is that \emph{sāyaṁ} is a worn-down \emph{sahāyaṁ}. It seems obvious. Obvious is not a citation. Chapter 31 drew the same line over a greeting: two sources for \textquotedblleft{}real and formal,\textquotedblright{} one for \textquotedblleft{}less common,\textquotedblright{} kept apart.
\end{culture}

\subsection*{Your turn}
Say these aloud:

\begin{itemize}
\raggedright
  \item \textquotedblleft{}nēnu mīku sahāyaṁ cēstānu\textquotedblright{} — I help you
  \item shorter, then polite — \textquotedblleft{}sāyaṁ cēstānu … sahāyaṁ cēyaṇḍi\textquotedblright{}
  \item the halves — \textquotedblleft{}saha … āya\textquotedblright{} — then \textquotedblleft{}exit … transit\textquotedblright{}
  \item two nouns, one verb — \textquotedblleft{}daya cēsi … sahāyaṁ cēyu\textquotedblright{}
  \item the sister that kept its own — \textquotedblleft{}udavi\textquotedblright{}
  \item the chapter — \textquotedblleft{}tīsukuṇṭānu, aḍugutānu, sahāyaṁ cēstānu\textquotedblright{}
\end{itemize}

\begin{tcolorbox}[breakable,colback=teal!4,colframe=teal!35!black,title={Before you move on}]
Say \textquotedblleft{}I help you,\textquotedblright{} name the case the helped person takes, and take \emph{sahāya} apart. (\emph{Nēnu mīku sahāyaṁ cēstānu}; the \textbf{dative -\te{కు}}; \emph{\textbf{\dv{सह}}} \textquotedblleft{}with\textquotedblright{} plus \emph{\textbf{\dv{आय}}} \textquotedblleft{}a going\textquotedblright{} — \textbf{exit} and \textbf{transit} carry the go-half.) Which other word is built on \emph{cēyu} that way, and which sister kept its own? (\emph{\textbf{\te{దయచేసి}}}; \textbf{Tamil}, with \emph{udavi}.) What may you \emph{not} say about \emph{sāyaṁ}? (\textbf{That it comes from \emph{sahāyaṁ}} — that derivation is not sourced here.)
\end{tcolorbox}
\section[nāku telugu iṣṭaṁ]{\te{నాకు} \te{తెలుగు} \te{ఇష్టం} — \textquotedblleft{}I like Telugu,\textquotedblright{} with no verb in it}

\label{lesson:TE-C34-ishtam}



\begin{quote}
\emph{Sahāyaṁ cēstānu} — a going with somebody. Say it, then \emph{tīsukuṇṭānu}. This last sentence has no verb.
\end{quote}

\subsection*{You'll want to know: The sentence}
\begin{quote}
\textbf{\te{నాకు} \te{తెలుగు} \te{ఇష్టం}.} \emph{nāku telugu iṣṭaṁ.} — \textquotedblleft{}\textbf{to me} — Telugu — \textbf{a liking}.\textquotedblright{}
\end{quote}

\textbf{\te{ఇష్టం}} is a \textbf{noun}: \textquotedblleft{}a liking, a wish,\textquotedblright{} and Telugu supplies no \emph{is}, having never needed one. Only the front piece moves — \emph{nāku}, \emph{nīku}, \emph{ataniki}. \textbf{\te{తెలుసు}} behaved so too: there the verb had no room for a person; here there is no verb.

\subsection*{The letters in this word}
\textbf{\te{ఇ}} is the standing vowel \emph{i}; \textbf{\te{ష్ట}} is the conjunct \textbf{\te{ష}} (\emph{ṣa}) over \textbf{\te{ట}} (\emph{ṭa}).

\begin{grammarlens}[title={Grammar Lens: the fourth dative sentence, and two ways round it}]
Four sentences, one idea: \emph{nāku telugu vaccu}, a \textbf{skill}; \emph{nāku telusu}, a \textbf{fact}; \emph{nāku arthamaindi}, a meaning that \textbf{landed}; \emph{nāku telugu iṣṭaṁ}, a \textbf{taste}. Chapter 6's dative subject was no curiosity: it is how Telugu files what happens \textbf{to} you.

A native verb does it too — \textbf{\te{నచ్చు}} (\emph{naccu}), \textquotedblleft{}to be pleasing\textquotedblright{}: \textbf{\te{నాకు} \te{ఇది} \te{నచ్చింది}} (\emph{nāku idi naccindi}). Tamil \textbf{\ta{நச்சு}} and Kannada \textbf{\ka{ನಚ್ಚು}} are the same verb: \emph{iṣṭaṁ} displaced nothing.

One shape takes an ordinary subject: \textbf{\te{ఇష్టపడు}} (\emph{iṣṭapaḍu}), \textbf{\te{నేను} \te{ఇష్టపడతాను}} — \emph{iṣṭaṁ} plus \textbf{\te{పడు}} (\emph{paḍu}), \textquotedblleft{}to fall,\textquotedblright{} Chapter 20's verb in \textbf{\te{వర్షం} \te{పడుతోంది}}. Liking is something you \textbf{fall into}.
\end{grammarlens}

\begin{cousinweb}
\textbf{\te{ఇష్టం}} is Sanskrit \textbf{\dv{इष्ट}} (\emph{iṣṭa}), \textquotedblleft{}desired,\textquotedblright{} from \textbf{\dv{इष्}} (\emph{iṣ}) — an Indo-European root whose English descendant is \textbf{ask}. So this chapter carried two asking-roots, not one: \emph{prach} behind \emph{pray}, \emph{iṣ} behind \emph{ask}. Telugu's \textbf{\te{అడుగు}} is cousin to neither: its relatives are Gondi and Kui.

Spanish \emph{me gusta}, Italian \emph{mi piace}, Hindi \emph{mujhe pasand hai}, Bengali \emph{āmār bhālo lāge}, Tamil \emph{eṉakkut tamiḻ piḍikkum} — and now \emph{nāku telugu iṣṭaṁ}. \textbf{Six languages, three unrelated families, one design}, with Chapter 6's Dravidian seam showing: \textbf{-ku}, \textbf{-ukku}, \textbf{-ge}, \textbf{-ikku}.
\end{cousinweb}

\subsection*{Your turn}
Say these aloud:

\begin{itemize}
\raggedright
  \item \textquotedblleft{}nāku telugu iṣṭaṁ\textquotedblright{} — \textquotedblleft{}to me, Telugu, a liking\textquotedblright{}
  \item three likers — \textquotedblleft{}nāku … nīku … ataniki telugu iṣṭaṁ\textquotedblright{}
  \item the four datives — \textquotedblleft{}vaccu … telusu … arthamaindi … iṣṭaṁ\textquotedblright{}
  \item native, then falling — \textquotedblleft{}naccindi … iṣṭapaḍatānu\textquotedblright{}
  \item the chapter — \textquotedblleft{}tīsukuṇṭānu, aḍugutānu, cēstānu, iṣṭaṁ\textquotedblright{}
  \item the endings — \textquotedblleft{}-ku · -ukku · -ge · -ikku\textquotedblright{}
\end{itemize}

\begin{tcolorbox}[breakable,colback=teal!4,colframe=teal!35!black,title={Before you move on}]
Say \textquotedblleft{}I like Telugu,\textquotedblright{} say how many verbs it holds, and name the four datives. (\emph{Nāku telugu iṣṭaṁ}; \textbf{none}; \emph{vaccu} a \textbf{skill}, \emph{telusu} a \textbf{fact}, \emph{arthamaindi} a \textbf{landing}, \emph{iṣṭaṁ} a \textbf{taste}.) Which native verb does the same job, and what hides inside \emph{iṣṭapaḍu}? (\emph{\textbf{\te{నచ్చు}}}; \emph{\textbf{\te{పడు}}}, the falling of \emph{varṣaṁ paḍutōndi}.) Which English word is \emph{iṣṭaṁ}'s cousin, and is it \emph{praśna}'s? (\textbf{Ask}; \textbf{no}, that is \textbf{pray} — and \textbf{six} languages in \textbf{three} families build liking so.)
\end{tcolorbox}
